\documentclass[11pt]{article}

\usepackage{amsmath,amssymb,graphicx}
\usepackage[margin=1in]{geometry}
\usepackage{booktabs}

\title{Phase II Boundary of Weak Reinforcement, Registration, and Delayed Memory on a Kernel-Induced Cochain Operator Architecture}
\author{Tomislav Rupi\'c \\ Independent Researcher}
\date{March 2026}

\begin{document}

\maketitle

\begin{abstract}
Phase II of the HAOS-IIP program was designed as a bounded search for the first persistence-supporting mechanism beyond the frozen pre-geometric limit of Phase I. The baseline architecture was kept fixed wherever possible: the same kernel-induced cochain support, the same projector convention, the same canonical representative packet families, and the same persistence-first promotion ladder. Only the law class was varied. Stage 19 tested weak unfrozen-kernel reinforcement across local scalar, mesoscopic basin-residence, transport-corridor, directional flux-driven, and mixed scalar-vector branches. Stage 20 formalized the resulting no-law boundary. Stage 21 tested selective interaction eligibility through phase-coherence, spatial-configuration, and composite-topology registration. Stage 22 tested delayed nonlocal kernel memory through single-lag retardation, exponential history, and finite-window averaging. Every family was numerically active and bounded. None produced a persistence gain. Across the executed pilots no composite lifetime increase appeared, no binding persistence increase appeared, no coarse basin-persistence increase appeared, no corridor dwell or locking increase appeared, no new stabilized ordering class appeared, and no refinement promotion was earned. The resulting statement is structural rather than anecdotal: within this kernel-induced cochain architecture, weak bosonic-style transport plus weak local, mesoscopic, directional, delayed, and registration-gated modifications do not generate persistence. Phase II therefore closes a broad weak-coupling mechanism family and justifies the move to Phase III as an operator-class change rather than another nearby feedback variant.
\end{abstract}

\section{Program Position}

Phase I closed with a sharp boundary: frozen pre-geometry can organize, but it does not bind. That result was strong because it survived a long sequence of diagnostics on the same frozen architecture, including weak effective binding extensions, relational feedback, phase synchronization, and minimal retarded memory. Phase II was the next lawful step. Instead of changing the ontology of the packets or the operator sector, it asked a smaller question first: if the architecture is allowed bounded weak modifications of substrate response, interaction eligibility, or temporal memory, does any persistence-supporting regime appear before a deeper architectural change becomes necessary?

The design discipline remained strict. The same three representatives were reused throughout late Phase II:
\begin{itemize}
\item clustered composite anchor,
\item phase-ordered symmetric triad,
\item counter-propagating corridor.
\end{itemize}
The same persistence-first observables were kept in force:
\begin{itemize}
\item composite lifetime,
\item binding persistence,
\item coarse basin persistence,
\item corridor dwell and transport locking,
\item emergence of new stabilized ordering classes.
\end{itemize}
The same promotion rule was kept: a branch earns refinement only if persistence improves before visual or secondary diagnostics are interpreted.

\section{Stage 19: Weak Substrate Reinforcement Family}

Stage 19 opened Phase II by changing the substrate law while preserving support, symmetry, projector convention, and representative set. Five weak reinforcement branches were executed.

\subsection{19A: Local Scalar Back-Reaction}

Stage 19A allowed the kernel weights to respond to a local energy-density signal through a slow bounded scalar reinforcement field. The branch was active but weak, with kernel deviation norms in the $10^{-6}$ to $10^{-5}$ range and maximum local deformation around $10^{-5}$. No persistence metric moved.

\subsection{19B-1: Basin-Residence Reinforcement}

Stage 19B Branch 1 replaced local occupation with mesoscopic basin-residence reinforcement. This was the strongest scalar-like branch in Phase II. Kernel deviation norms rose into the $10^{-3}$ range and maximum local deformation reached approximately $5\times10^{-2}$. Even at that scale, composite lifetime, binding persistence, and coarse basin persistence remained unchanged.

\subsection{19B-2: Transport-Corridor Reinforcement}

Stage 19B Branch 2 reinforced recently used transport pathways through a smoothed flux-magnitude memory field. The corridor field was alive, but the deformation scale fell back into the $10^{-6}$ to $10^{-5}$ range. No corridor-selective locking appeared, and the persistence ladder remained flat.

\subsection{19C: Directional Flux-Driven Reinforcement}

Stage 19C changed symmetry class and introduced a bounded directional substrate response. The resulting directional field and deformation were measurable but weak, with kernel deviation norms again in the $10^{-6}$ range and local deformation in the $10^{-5}$ range. A small deflection proxy shift appeared in the corridor case, but no dwell gain, no transport locking, and no persistence improvement followed.

\subsection{19D: Mixed Scalar-Vector Reinforcement}

Stage 19D combined scalar density memory and vector-like flux memory in a single bounded update law. This branch entered the substrate strongly. Mixed field norms rose into the $10^{-3}$ to $10^{-2}$ range, kernel deviation norms reached approximately $9.4\times10^{-3}$, and maximum local deformation reached roughly $3\times10^{-2}$. The persistence outcome remained null.

\subsection{Stage 19 Boundary}

Stage 19 is not a dead-end null. It shows a stronger result. Weak substrate-law reinforcement can enter the architecture, produce bounded nonzero deformation, and still fail to create historical stabilization. The negative therefore does not depend on a single symmetry class. It survives local scalar reinforcement, mesoscopic scalar reinforcement, directional flux reinforcement, and mixed scalar-vector reinforcement.

\section{Stage 20: No-Law Boundary Consolidation}

Stage 20 did not add a new mechanism. It formalized the boundary exposed by Stage 19 and recorded it as a structural result. The report makes four points that matter for the rest of the program:
\begin{enumerate}
\item the null is regime-level, not merely visual;
\item bounded substrate deformation is real across multiple branches;
\item persistence does not follow from deformation in the tested class;
\item further nearby scalar-like reinforcement variants are low-priority.
\end{enumerate}

This matters because it converts what could have looked like a sequence of small nulls into a mapped forbidden region of mechanism space.

\section{Stage 21: Mutual Registration and Constraint-Coupled Interaction}

Stage 21 left reinforcement space and tested a different question. Perhaps persistence does not require a stronger medium at all. Perhaps it requires selective interaction eligibility: interaction becomes active only when packet structures satisfy a compatibility condition.

Three registration classes were tested:
\begin{itemize}
\item phase-coherence registration,
\item spatial-configuration registration,
\item composite-topology registration.
\end{itemize}
The law class remained bounded and memoryless:
\[
F_{AB}(t)=\varepsilon\,R_{\mathrm{align}}(A,B,t)\,F_{\mathrm{base}}(A,B,t),
\qquad \varepsilon\in\{0.01,0.02\}.
\]

The 27-run base pilot remained fully negative at the persistence layer. No composite lifetime gain, no binding persistence gain, no coarse basin gain, and no corridor dwell gain appeared. No stabilized ordering class appeared. No refinement promotion was earned.

The negative is nevertheless informative because the gates were active. Phase-coherence registration reached bounded mean strengths up to $0.1829$ but never opened a sustained duty cycle. Spatial-configuration registration opened full eligibility windows with duty cycle $1.0$ and effective interaction-energy fractions up to $1.34\times10^{-2}$, yet still produced no persistence gain. Composite-topology registration also opened nontrivial windows and still failed to move the regime-level observables. Stage 21 therefore sharpens the architecture picture: the system can detect compatibility without converting it into capture.

\section{Stage 22: Delayed Nonlocal Kernel Memory}

Stage 22 asked the strongest weak-memory question yet tested in the program. If persistence does not emerge from instantaneous weak reinforcement or instantaneous weak registration, can a temporally accumulated nonlocal memory field create the first anchoring channel?

Three delayed-memory laws were executed:
\begin{itemize}
\item single-lag retarded reinforcement,
\item exponential-history kernel memory,
\item finite-window historical averaging.
\end{itemize}
All three used the same smoothed nonlocal energy-density driver and the same weak coupling ladder.

The result remained uniformly negative at the persistence level. Across all 27 runs:
\begin{itemize}
\item composite lifetime delta stayed $0.0000$,
\item binding persistence delta stayed $0.0000$,
\item coarse basin persistence delta stayed $0.0000$,
\item corridor dwell delta stayed $0.0000$,
\item no new stabilized ordering class appeared,
\item no refinement promotion was earned.
\end{itemize}

Again, the null is not an under-activation artifact. The single-lag and finite-window classes reached memory-field norms near $1.35\times10^{-1}$ and kernel deviation norms near $2.7\times10^{-3}$. The exponential-history class was stronger, reaching memory-field norm $2.65\times10^{-1}$, kernel deviation norm $5.3\times10^{-3}$, and local deformation $1.52\times10^{-2}$. None of these active delayed-memory fields altered persistence.

\section{Phase II Comparison}

The strongest summary of Phase II is not that one branch failed. It is that several physically natural weak-law families were tried under a consistent measurement ladder and all stayed persistence-null.

\begin{table}[h]
\centering
\small
\begin{tabular}{@{}p{1.6in}p{1.7in}p{1.6in}p{1.0in}@{}}
\toprule
Stage / branch & New ingredient & Strongest active scale & Persistence outcome \\
\midrule
19A & local scalar reinforcement & kernel deviation $\sim 10^{-6}$ to $10^{-5}$ & negative \\
19B-1 & basin-residence reinforcement & kernel deviation $\sim 10^{-3}$ to $10^{-2}$ & negative \\
19B-2 & transport-corridor reinforcement & kernel deviation $\sim 10^{-6}$ to $10^{-5}$ & negative \\
19C & directional flux reinforcement & kernel deviation $\sim 10^{-6}$ & negative \\
19D & mixed scalar-vector reinforcement & kernel deviation $\sim 10^{-2}$ & negative \\
21 phase & phase registration gating & interaction fraction $\sim 3.3\times10^{-3}$ & negative \\
21 spatial & spatial registration gating & duty cycle $1.0$, interaction fraction $\sim 1.3\times10^{-2}$ & negative \\
21 topology & motif registration gating & duty cycle $1.0$, interaction fraction $\sim 7.3\times10^{-3}$ & negative \\
22A & single-lag delayed memory & kernel deviation $\sim 2.7\times10^{-3}$ & negative \\
22B & exponential-history memory & kernel deviation $\sim 5.3\times10^{-3}$ & negative \\
22C & finite-window delayed memory & kernel deviation $\sim 2.7\times10^{-3}$ & negative \\
\bottomrule
\end{tabular}
\caption{Phase II mechanism summary. All branches were bounded and numerically active. None produced a gain in composite lifetime, binding persistence, coarse basin persistence, or corridor locking.}
\end{table}

This table is the key Phase II result. The architecture is responsive. It can carry local deformation, mesoscopic deformation, directional bias, gating windows, and delayed memory. What it does not do is turn any of those into persistence.

\section{What Phase II Falsifies}

Phase II does not prove that every classical or continuous-wave model must fail. That claim would exceed the evidence. The falsified statement is narrower and stronger:

\begin{quote}
Within this kernel-induced cochain architecture, weak bosonic-style transport plus weak local, mesoscopic, directional, delayed, and registration-gated modifications do not generate persistence.
\end{quote}

That means the following families are now low-priority in their weak bounded forms:
\begin{itemize}
\item nearby scalar reinforcement variants,
\item nearby directional reinforcement variants,
\item nearby weak gating variants,
\item nearby weak delayed-memory variants.
\end{itemize}

This is why the Phase II negative plateau is valuable. It is not empty. It removes a broad class of nearby explanations and protects the program from coefficient drift.

\section{Why Phase III Is Justified}

The next coherent step is not another weak reinforcement law. It is an operator-class change. The strongest motivation is not that a Dirac--K\"ahler sector must work, but that the existing packet ontology and transport logic have now been tested across a broad family of weak modifications without producing persistence.

Phase III is therefore justified as a new baseline question:

\begin{quote}
If the same canonical packet families are propagated under a validated Dirac--K\"ahler sector, does the persistence landscape change qualitatively before any later binding-specific mechanism is introduced?
\end{quote}

This is the cleanest Phase III opener because it changes the operator class while keeping the rest of the comparative apparatus intact. It does not claim fermionic binding. It asks whether first-order graded propagation changes morphology, overlap persistence, or escape behavior at all.

\section{Conclusion}

Phase II closes as a coherent boundary program. It explored the largest nearby family of weak architectural modifications that could plausibly have created persistence without changing the operator sector itself. The result is consistent across every tested branch:
\begin{itemize}
\item the branches are active,
\item the deformations and gates are bounded,
\item the architecture responds,
\item persistence does not follow.
\end{itemize}

The strongest compact statement is:

\begin{quote}
Weak substrate reinforcement, selective interaction eligibility, and delayed nonlocal memory are all insufficient, at the tested couplings, to generate persistence in the current kernel-induced cochain architecture.
\end{quote}

That is not a failure to find the right coefficient. It is a mapped boundary. Phase III is therefore warranted as a change of operator class rather than another nearby weak-law search.

\end{document}
